% STATUS: DRAFT
\chapter{Classification: types I, II, III}\label{ch:classification}

\begin{physmotiv}
The question is: \emph{how many fundamentally different kinds of quantum systems are there?} Answer (Murray and von Neumann, 1936--43): factors fall into three types, distinguished by which \emph{dimensions} projections can have. Type I is textbook QM. Type II has a continuous dimension and a trace but no minimal observables. Type III has neither trace nor dimension: every non-trivial question has ``infinite dimension'', and entropy is meaningless. Local algebras in QFT are type III, and the program of this book is to understand how gravity turns III into II.
\end{physmotiv}

\section{Comparing projections}

Two projections $e,f\in\M$ are \emph{Murray--von Neumann equivalent}, $e\sim f$, if there is a partial isometry $v\in\M$ with $v\ad v=e$ and $vv\ad=f$. That is, $\M$ itself contains an operator that maps the range of $e$ isometrically onto the range of $f$. Write $e\preceq f$ if $e\sim f_0\le f$. This is \emph{dimension-comparison inside $\M$}: in $\BH$ it says $\mathrm{rank}\,e=\mathrm{rank}\,f$, but in a smaller algebra it need not, because the required partial isometry might not belong to $\M$.

\begin{theorem}[comparison]
In a factor, any two projections are comparable: $e\preceq f$ or $f\preceq e$.
\end{theorem}
So the projections of a factor are totally ordered by size; this ordering is captured by a \emph{dimension function} $d:P(\M)\to[0,\infty]$, unique up to scale, with $d(e)=d(f)$ iff $e\sim f$, and additive on orthogonal projections. A projection is \emph{finite} if it is not equivalent to a proper subprojection of itself, i.e.\ not a ``Hilbert-hotel'' projection (\cref{ch:hilbert}). It is \emph{minimal} if it dominates no non-zero smaller projection.

\section{The three types}

\begin{center}
\renewcommand{\arraystretch}{1.35}
\begin{tabular}{@{}l p{3.8cm} l l@{}}
\toprule
Type & Definition & Range of $d$ & Example\\
\midrule
I$_n$ & has minimal projections; $\id$ finite & $\{0,1,\dots,n\}$ & $M_n(\C)$\\
I$_\infty$ & has minimal projections; $\id$ infinite & $\{0,1,2,\dots,\infty\}$ & $\BH$, $\dim\HH=\infty$\\
II$_1$ & no minimal projections; $\id$ finite & $[0,1]$ & hyperfinite II$_1$ factor $R$\\
II$_\infty$ & no minimal projections; $\id$ infinite, but finite projections exist & $[0,\infty]$ & $R\bar\otimes\BH$\\
III & every non-zero projection is infinite (and all are equivalent to $\id$) & $\{0,\infty\}$ & QFT, Powers, Araki--Woods\\
\bottomrule
\end{tabular}
\end{center}

\begin{itemize}
\item \textbf{Type I:} $\M\cong\Bnd{K}$ for a Hilbert space $K$; i.e.\ $\M=\BH_A\otimes\id_B$ on $\HH=\HH_A\otimes\HH_B$. This is the world of tensor factorization, density matrices, and $\Tr$.
\item \textbf{Type II:} a continuous dimension. In II$_1$ there are projections of ``dimension $1/\pi$''; there is a unique tracial state $\tau$ with $\tau(e)=d(e)$. There is no pure state. The trace is finite on $\id$ and tracial states are maximally mixed states.
\item \textbf{Type III:} no trace (not even on non-zero projections), $d$ takes only the values $0,\infty$. Every non-zero projection $e$ is equivalent to $\id$: any non-trivial yes/no question is as big as the entire system.
\end{itemize}

\begin{whatwrong}
In type III, the statement ``$e\sim\id$ for all non-zero $e$'' means we can \emph{locally} (inside the algebra) map the entire space into the range of an arbitrarily small projection. A tiny region of QFT contains operators that can prepare, up to unitary equivalence, anything. This is the algebraic source of the Reeh--Schlieder property and of the divergence of entanglement entropy.
\end{whatwrong}

\section{Where the types come from: the hyperfinite examples}

\textbf{II$_1$.} Take the infinite spin chain $M_{2^\infty}$ with the tracial state $\tau$ (\cref{ch:gns}). The weak closure $R=\pi_\tau(\A)''$ of its GNS representation is the \emph{hyperfinite II$_1$ factor}: ``hyperfinite'' means it is the weak closure of an increasing union of finite-dimensional algebras; it is unique (Murray--von Neumann, Connes). It has a unique tracial state, and dimension values $d(e)\in[0,1]$: a projection of dimension $k/2^n$ exists inside $M_{2^n}$ (rank $k$), and these are dense in $[0,1]$. Physically $R$ is the algebra of observables of the infinite chain at \emph{infinite temperature}.

\textbf{II$_\infty$.} $R\bar\otimes\BH$ with trace $\tau\otimes\Tr$: infinite temperature plus an infinite ordinary system. In \cref{ch:takesaki} the crossed product of a type III$_1$ factor by its modular flow is of this form.

\textbf{III.} For $p\ne1/2$ the state $\omega_p=\bigotimes_j\Tr(\mathrm{diag}(p,1-p)\,\cdot)$ has GNS von Neumann algebra the \emph{Powers factor} $R_\lambda$ with $\lambda=\min(p,1-p)/\max(p,1-p)\in(0,1)$, of type III$_\lambda$ (\cref{ch:connes_class}). Intuitively: one tries to renormalize the finite-volume trace $\Tr(\rho_N\,\cdot)\to\Tr(\cdot)$ but the ratios $\rho_N/\Tr\rho_N$ spread over exponentially many scales, so that no consistent trace survives. For example, the non-interacting chain $H=\frac h2\sum\sigma^z_j$ at any finite temperature gives such a product state with $\lambda=\ee^{-\beta h}$, hence type III$_\lambda$ in infinite volume.

\begin{keyidea}
Type is decided by what dimensions projections can have: I: integers; II$_1$: $[0,1]$; II$_\infty$: $[0,\infty]$; III: $\{0,\infty\}$. Type is intrinsic to the algebra, unchanged by choosing different faithful representations. Local QFT algebras are type III, so no trace and no entropy.
\end{keyidea}

\section{What ``entropy exists'' means for each type}

\begin{center}
\renewcommand{\arraystretch}{1.3}
\begin{tabular}{@{}llp{6.5cm}@{}}
\toprule
Type & Trace & Entropy $S(\omega)=-\tau(\rho_\omega\log\rho_\omega)$\\
\midrule
I$_n$ & finite & yes; $0\le S\le\log n$; maximum at $\rho=\id/n$\\
I$_\infty$ & semifinite & yes for $\rho$ trace-class with finite entropy; unbounded above\\
II$_1$ & finite, $\tau(\id)=1$ & yes; $S\le0$; maximum $0$ at the tracial state\\
II$_\infty$ & semifinite & yes for $\rho$ in the domain; unbounded in both directions, state-independent shift ambiguity from the scale of $\tau$\\
III & none & \emph{no density matrix exists}; relative entropy still exists (\cref{ch:modham})\\
\bottomrule
\end{tabular}
\end{center}
The sign: in II$_1$ with normalized trace $S\le0$ is a \emph{relative} entropy with respect to the maximally mixed state, shifted. The gravitational examples will have $S=\Sgen+\text{const}$ with a maximum at the de Sitter vacuum (\cref{ch:clpw}).

\section{Non-factors and the general case}

A general von Neumann algebra decomposes as a direct integral of factors (\cref{ch:factors}), and a factor has exactly one type. A general algebra thus has a decomposition $\M=\M_I\oplus\M_{II}\oplus\M_{III}$ into homogeneous pieces of the respective types, with possibly types I$_n$ for different $n$ in the type I part. For physical systems with superselection sectors, the sectors can have different types.

\section*{Exercises}
\begin{exercise}
Show that in $\BH$ ($\HH$ separable infinite-dimensional) two projections are equivalent iff they have the same rank (possibly $\infty$). Show the projection onto the even sites of $\ell^2$ is equivalent to $\id$ but not equal to it. Is it finite?
\end{exercise}
\begin{exercise}
In $R=\pi_\tau(M_{2^\infty})''$ show that $\tau(e)=k/2^n$ for a rank-$k$ projection in $M_{2^n}$, and use density to argue that $d$ takes all values in $[0,1]$ (every real $r\in[0,1]$ is the dimension of some projection). Why does this prove that there is no minimal projection?
\end{exercise}
\begin{exercise}
Show that a type I$_n$ factor has a unique tracial state, and compute the entropy of $\omega=\tau(\rho\,\cdot)$ when $\rho=\mathrm{diag}(p_1,\dots,p_n)\cdot n$. (Note the normalization $\tau=\Tr/n$ shifts $S$ by $\log n$ compared to Shannon.)
\end{exercise}
\begin{exercise}
Argue that if $\M$ is type III then no pure normal state exists and no trace exists. (Pure normal state $\Rightarrow$ minimal projection; trace $\Rightarrow$ finite projections.)
\end{exercise}
\begin{exercise}
Why can the algebra of a bounded region of QFT not be type I? (Argue using the existence of minimal projections ($\Rightarrow$ a pure normal state) versus the fact that the vacuum is cyclic and separating and entangled; see \cref{ch:typeIII_local}.)
\end{exercise}
